\chapter{The Tiered Proof Pipeline: Methods and Lessons}\label{ch:pipeline}

\Cref{ch:architecture} described the gates a theorem must pass before it carries the badge
\TA. This chapter answers the question that comes before the gates: \emph{how were the 99
theorems of the library \lean{DualScaleStream2} actually produced, by whom, at what cost,
and what went wrong on the way?} The answer is a method, not an anecdote. Between
2026-09-16 21:27 and 2026-09-17 05:07 UTC a single orchestrating session wrote the
statements, checked their conventions in a computer-algebra system, and handed the proofs
to three tiers of language models---a 7-billion-parameter prover running on a local GPU, a
cheap hosted model, and a more capable hosted model---escalating a goal to the next tier
only after the cheaper one had failed. The Lean kernel was the only judge at every step.
The run was instrumented: every attempt by the local prover was logged, every report from a
hosted model was recompiled, and the per-tier results were written down as measured, not as
expected \source{STREAM2\allowbreak\_WORKFLOW.md §0, ll.~3--5}. A student should care because this
is the first place in the book where the \emph{production} of formal mathematics is treated
as an experiment with controls, and because each of the safeguards now in the repository
exists as the answer to a specific, recorded failure. The chapter is written so that a
student with a Lean project, a GPU and an account with a model provider can repeat the
procedure on a library of their own.

%=====================================================================================
\section{Many small goals, one judge}\label{sec:pipeline:judge}

\subsection{What the pipeline is for}

The library \lean{DualScaleStream2}\index{DualScaleStream2@\lean{DualScaleStream2}} was
planned as ``a kernel-checked Lean 4 layer for the \emph{mathematical skeleton} of type II /
heterotic string theory on $\KT$---charge lattices, their duality groups, the doubled (DFT)
algebra, flux/tadpole counting, and the Mathieu-moonshine numerics'' \source{STREAM2\allowbreak\_WORKFLOW.md
§1, ll.~9--13}. The same paragraph fixes what it is not: ``Not a goal, and never to be
claimed: `formalizing string theory' as physics'' \source{STREAM2\allowbreak\_WORKFLOW.md §1,
ll.~15--18}. The plan lists six phases P2.1--P2.6 (lattices, $T^2$ dualities, K3 moduli,
DFT algebra, flux and tadpole, moonshine interface), each pinned to a paper section
\source{STREAM2\allowbreak\_WORKFLOW.md §5, ll.~67--76}, and a checklist of eighteen roadmap items whose
``Done'' means ``every theorem for the item is committed, sorry-free and axiom-audited''
\source{STREAM2\allowbreak\_WORKFLOW.md §5.1, ll.~78--103}.

The work therefore consisted of a large number of small, independent goals: 99 designed
theorems at the final gate \source{STREAM2\allowbreak\_WORKFLOW.md §6, ll.~185--189}. Most are the
kind of statement a student meets in the earlier chapters of this book---that a fixed
$8\times8$ integer matrix has a given inverse, that $\eta^2=1$, that a block-matrix product
equals another block matrix, that a table of 26 integers has a given sum of squares. A few
need an idea. The economic problem is to close the many cheap goals cheaply and spend
effort only on the few that need it, without ever letting a wrong proof in.

\subsection{The hard constraints}\label{sec:pipeline:constraints}

Seven rules were fixed before the run, each traced to an incident in the repository's own
history \source{STREAM2\allowbreak\_WORKFLOW.md §2, ll.~20--38}. Four matter for this chapter.
\begin{enumerate}[leftmargin=*,itemsep=2pt]
\item \textbf{Kernel is the only accept gate.} ``A goal is closed when \texttt{lake build
  DualScale\allowbreak Stream2} is green, the zero-\lean{sorry} audit passes, and
  \texttt{\#print axioms} shows only \lean{propext}, \lean{Classical.choice},
  \lean{Quot.sound}. No \lean{native\_decide}'' (ll.~22--24).
\item \textbf{Producer $\ne$ verifier.}\index{producer $\ne$ verifier} ``No agent or
  model grades its own output; the gate runs separately'' (ll.~25--26).
\item \textbf{No citation from memory.} Every Tier L docstring cites a file in
  \texttt{papers/foundations/} and a line range a script can \texttt{grep} (ll.~27--28).
\item \textbf{Separate library.} Stream 2 lives in its own \lean{lean\_lib}, imports Stream
  1 and never the reverse, so ``Stream 1's green core cannot be turned red by Stream 2
  work'' (ll.~32--33).
\end{enumerate}
The first rule is the one to keep in mind throughout. A language model's proof is a
\emph{candidate}; a candidate becomes a theorem at the moment the kernel accepts it, and at
no other moment. Everything else in the pipeline---tiers, routing, prompts, budgets---is
about generating candidates cheaply. Nothing in it is about deciding correctness.

\subsection{The four tiers}\label{sec:pipeline:tiers}

The workflow assigns every task to one of four model tiers \source{STREAM2\allowbreak\_WORKFLOW.md §3,
ll.~40--50}; \cref{tab:pipeline:tiers} reproduces the assignment. The names T0--T3 are this
project's and denote cost, not intelligence: T3 is free once the GPU is bought, T0 is the
most expensive model and is ``kept to a few calls per package''.

\begin{table}[ht]\centering\small
\caption{Model tiers of the Stream 2 workflow \source{STREAM2\allowbreak\_WORKFLOW.md §3, ll.~42--50}.
The local provers ran under Ollama on an NVIDIA T4; the quantization suffixes (Q8\_0, Q6\_K)
denote the 8-bit and 6-bit weight formats used.}
\label{tab:pipeline:tiers}
\begin{tabular}{@{}p{0.6cm}p{3.3cm}p{7.9cm}@{}}\toprule
Tier & Model & Used for \\ \midrule
T0 & Opus (orchestrator) & selecting the work package; the \emph{statement review} (``does
the Lean say what the paper says?''); tier labels A/L/C; commit \\[2pt]
T1 & Sonnet & proof bodies needing a real idea (block-matrix algebra, involution arguments);
residual goals after T3/T2 \\[2pt]
T2 & Haiku (default worker) & statement extraction from pinned paper lines into
\lean{sorry}-stubs; residual \emph{mechanical} goals; paper fetch, citation-line checks \\[2pt]
T3 & DeepSeek-\allowbreak Prover-V2-7B Q8\_0, Goedel-\allowbreak Prover-V2-8B Q6\_K
(local) & first pass on every open goal via \texttt{tools/prover\_loop.py} \\
\bottomrule\end{tabular}
\end{table}

The default routing is ``T3 $\to$ T2 $\to$ T1, escalating a goal only after the cheaper tier
has failed under the kernel gate. T0 never writes proofs in the steady state''
\source{STREAM2\allowbreak\_WORKFLOW.md §3, ll.~49--50}. \Cref{fig:pipeline:flow} draws the loop. The
picture to keep is that of a filter cascade: each tier removes the goals it can close, and
only the residue reaches the next, more expensive stage. Whether the cascade is efficient
is an empirical question, answered in \cref{sec:pipeline:results}.

\begin{figure}[ht]\centering
\begin{tikzpicture}[>=Stealth, node distance=7mm and 9mm,
  box/.style={draw, rounded corners=1.5pt, minimum height=8mm, align=center, font=\small, inner sep=3pt},
  gate/.style={draw, fill=green!8, diamond, aspect=2.2, align=center, font=\footnotesize, inner sep=1pt},
  lab/.style={font=\footnotesize, align=center}]
  \node[box, fill=blue!6] (stmt) {statements with \lean{sorry}\\ (T2 extracts, T0 reviews, sympy checks conventions)};
  \node[box, below=of stmt] (t3) {T3 local prover\\ \texttt{prover\_loop.py}};
  \node[gate, right=of t3] (k3) {kernel};
  \node[box, below=of t3] (t2) {T2 Haiku};
  \node[gate, right=of t2] (k2) {kernel};
  \node[box, below=of t2] (t1) {T1 Sonnet};
  \node[gate, right=of t1] (k1) {kernel};
  \node[box, fill=green!8, right=22mm of k2] (done) {closed goal\\ G3: build, \lean{sorry} audit,\\ axiom audit, statement lock};
  \draw[->] (stmt) -- (t3);
  \draw[->] (t3) -- (k3);
  \draw[->] (k3) -- node[lab, above, sloped] {accepted} (done.west |- k3) -- (done.west |- k3);
  \draw[->] (k3) -- node[lab, right] {rejected} (k3 |- t2.north) -- (t2.north east);
  \draw[->] (t2) -- (k2);
  \draw[->] (k2) -- node[lab, above] {accepted} (done);
  \draw[->] (k2) -- node[lab, right] {rejected} (k2 |- t1.north) -- (t1.north east);
  \draw[->] (t1) -- (k1);
  \draw[->] (k1) -- node[lab, above, sloped] {accepted} (done.west |- k1);
  \draw[->, dashed] (k1) -- ++(0,-9mm) node[lab, below] {residue: orchestrator re-derives the route (\cref{sec:pipeline:route})};
\end{tikzpicture}
\caption{The tiered cascade. Every arrow into ``closed goal'' passes through a kernel compile
that the producing tier does not run; a rejected candidate sends the \emph{unchanged}
statement, with its proof restored to exactly \lean{sorry}, to the next tier.}
\label{fig:pipeline:flow}
\end{figure}

%=====================================================================================
\section{Statement design before proof search}\label{sec:pipeline:design}

The per-package pipeline has nine steps S0--S9 \source{STREAM2\allowbreak\_WORKFLOW.md §4, ll.~52--65}.
Steps S0--S3 happen before any proving: select one paper section for one module (S0, T0);
pin the source lines (S1, T2); extract Lean statements with \lean{sorry} bodies and
tier-labelled docstrings, gate G1 being that the file ``builds with only \lean{sorry}
warnings'' (S2, T2); and review each statement, gate G2 being ``recorded sign-off (the one
step no kernel can check: mis-formalization)'' (S3, T0). Only then come the proof search
steps S4--S6, verification S7, recording S8 and commit S9.

The lesson list makes the ordering a rule for the next run: ``Statement design first, in one
pass. Write every statement with sympy-checked conventions and \lean{sorry}, then run
\texttt{lake build}, \texttt{statement\_lock --update}, and T0 review. Only after that does any
proving start. This run interleaved design and proving, which caused re-locks''
\source{LL.md §S2-I.1, ll.~118--121}. Two design practices deserve a closer look.

\subsection{Pre-verify the convention symbolically}\label{sec:pipeline:sympy}

A theorem about block matrices depends on sign and ordering conventions: where the $B$-field
sits in the generalized metric, whether a group element acts as $gHg^{\mathsf T}$ or
$g^{\mathsf T}Hg$, which way a shift is oriented. A statement with the wrong convention is
\emph{false}, and a prover set on a false statement fails in a way that is
indistinguishable, from the outside, from a prover that is not clever enough. The log puts
the cost plainly: ``A sign slip caught this way costs minutes. The same slip caught only
after an agent has spent an hour on an unprovable statement costs far more''
\source{LL.md §S2.12, ll.~101--105}. So ``every block-matrix convention \ldots\ was checked
in sympy before its Lean statement was written'' (ll.~102--103).

\begin{example}[The $B$-shift convention, checked before it was stated]\label{ex:pipeline:bshift}
The module \texttt{DFT/\allowbreak BShift.\allowbreak lean} proves that the integer $B$-field shift acts on the
generalized metric by the $\Theta$-shift. Its header records that ``of the four candidates
$gHg^{\mathsf T}$/$g^{\mathsf T}Hg$ with $\pm\Theta$ in the upper or lower block, exactly this
one holds, so the sign/ordering here is not a guess'' \source{BShift.lean, ll.~11--13}. The
author of this chapter repeated the check. With the book's convention (\cref{ch:genmetric})
\[
\HH(G,B)=\begin{pmatrix} G-BG^{-1}B & BG^{-1}\\ -G^{-1}B & G^{-1}\end{pmatrix},\qquad
g_\Theta=\begin{pmatrix} 1 & \Theta\\ 0 & 1\end{pmatrix},
\]
take $d=2$, a generic symmetric $G=\begin{psmallmatrix} g_{11}&g_{12}\\ g_{12}&g_{22}
\end{psmallmatrix}$ and generic antisymmetric $B=\begin{psmallmatrix}0&b\\-b&0
\end{psmallmatrix}$, $\Theta=\begin{psmallmatrix}0&t\\-t&0\end{psmallmatrix}$, and ask sympy
which of the four expressions equals $\HH(G,B+\Theta)$ identically in $(g_{11},g_{12},
g_{22},b,t)$. The result is
\[
g_\Theta\,\HH(G,B)\,g_\Theta^{\mathsf T}=\HH(G,B+\Theta)\ \text{(identity)},\qquad
g_\Theta^{\mathsf T}\HH g_\Theta,\ g_{-\Theta}\HH g_{-\Theta}^{\mathsf T},\
g_{-\Theta}^{\mathsf T}\HH g_{-\Theta}\ \text{(all fail)}.
\]
The surviving candidate is the statement of \lean{genMetric\_bshift} \source{BShift.lean,
ll.~76--78}, and the same script confirms $\eta\HH\eta\HH=1$ and $g_\Theta^{\mathsf T}\eta
g_\Theta=\eta$ for antisymmetric $\Theta$ but not for symmetric $\Theta$. Four lines of
computer algebra rule out three false theorems before any prover sees them. The check is a
\emph{symbolic check}: it is evidence for the statement's truth at $d=2$, not a proof for
general $d$, and the log insists that ``papers must label it as a symbolic check and must
not quote it as a theorem'' \source{LL.md §S2.12, ll.~105--107}. The proof for every $d$ is
the kernel's, \cref{sec:pipeline:machine}.
\end{example}

\subsection{Split the concrete from the symbolic}\label{sec:pipeline:split}

The second design practice was learned rather than planned. After the first phase the
orchestrator observed that the local prover ``closed essentially every \emph{concrete} goal
(fixed matrices, numerals, \lean{fin\_cases} tables, the 64-entry E8 inverse check) and
essentially none of the \emph{symbolic} general-$n$/general-$d$ goals'', and drew the rule:
``write statements so the concrete content is split into separate lemmas---they are nearly
free at T3'' \source{LL.md §S2.1, ll.~25--31}. A theorem that says ``the matrix
\lean{cartanE8} is unimodular'' is symbolic if stated through the determinant and concrete
if stated as ``\lean{cartanE8Inv * cartanE8 = 1}'' for an explicit certificate matrix; the
general lemma \lean{isUnimodular\_of\_mul\_eq\_one} then turns the concrete fact into the
abstract one once, for every lattice. The certificate itself was computed by sympy and
``carries no mathematical guarantee on its own'' \source{E8.lean, ll.~105--108}; the
64-entry check is what the kernel does.

\begin{pitfallbox}[A provable statement is not yet the right statement]
The kernel checks a proof against a statement. It does not check the statement against the
paper, and it does not check that the statement is still the one that was reviewed. Gate G2
(a reader) and the statement lock (\texttt{tools/statement\_lock.py}, \cref{ch:architecture})
cover the two gaps. The run recorded both failure modes: a docstring that drifted from
``if'' to ``iff'' \source{LL.md §S2.13, ll.~109--112}, and re-locks caused by proving before
the design was frozen \source{LL.md §S2-I.1, ll.~118--121}. Exercise~\ref{ex:pipeline:iff}
asks the student to find a surviving instance of the first kind in the current source.
\end{pitfallbox}

%=====================================================================================
\section{The local prover loop}\label{sec:pipeline:loop}

Step S4 is a program, \texttt{tools/prover\_loop.py}, 188 lines of Python that a student can
read in full \source{tools/prover\_loop.py, ll.~1--188}. Its docstring states the contract:
for every \lean{theorem}/\lean{lemma} whose body is \lean{sorry}, ``ask one or more local
Ollama prover models for a proof, splice each candidate in, and compile with \lean{lake env
lean}. A candidate is ACCEPTED only if the file compiles with zero errors and exactly one
fewer \lean{sorry} warning than before. The model never grades itself; the Lean kernel is the
only judge'' (ll.~5--9). We go through the loop in the order the program executes it,
because each clause encodes a decision.

\paragraph{Finding the goals.} A regular expression matches a declaration head from
\lean{theorem} or \lean{lemma} through \lean{:= by} followed by an indented \lean{sorry}
(ll.~39--44). The pattern is written so that the statement ``must not run into the next
declaration---otherwise, once an earlier goal is solved, its match stretches forward and
swallows the next goal's \lean{sorry} (observed in the first smoke test: 2 of 7 goals
silently skipped)'' (ll.~36--38). This is the first of several places where the tool's
comments are themselves a lab notebook.

\paragraph{The baseline.} Before searching, the file is compiled once to record its error
count and its number of \lean{sorry} warnings; a file that does not compile is skipped
outright (ll.~133--136). The accept criterion will be stated relative to this baseline.

\paragraph{The prompt.} The model receives the file text up to and including the
\lean{:= by} of the current goal, framed by one fixed instruction: ``Complete the following
Lean 4 code (Lean 4 + Mathlib). Replace the final \lean{sorry} with a complete tactic proof.
Do not use \lean{sorry}, \texttt{admit}, or \lean{native\_decide}. Output the complete final
theorem in one \texttt{lean4} code block'' (ll.~70--75). The context window is 8192 tokens
and at most 2048 tokens are generated by default (ll.~80, 124). There is no mathematical
hint in the prompt: the T3 tier sees the statement and the file's preceding definitions,
nothing else.

\paragraph{Extraction and the forbidden list.} The reply is scanned for code blocks, and the
last block containing a declaration with the goal's name is taken; its body after
\lean{:= by} up to the next declaration keyword is the candidate proof (ll.~89--100). A
candidate containing \lean{sorry}, \texttt{admit}, \lean{native\_decide} or an \texttt{axiom}
line is rejected before compilation (l.~35, ll.~162--164). The reason for the last two is
the axiom audit of \cref{ch:architecture}: \lean{native\_decide} would pass the build and
fail the audit, and an \texttt{axiom} would make anything provable.

\paragraph{The accept criterion.} The candidate is spliced in at the indentation of the
original \lean{sorry} (ll.~103--110) and the whole file is compiled. It is accepted if and
only if the error count is zero \emph{and} the \lean{sorry}-warning count is exactly the
baseline minus one (l.~167). The second clause matters: a candidate that ``proves'' the goal
by leaving a nested \lean{sorry}, or that accidentally deletes a neighbouring goal, changes
the count by the wrong amount and is rejected even though the file compiles. On acceptance
the file is written back and the new count becomes the baseline (ll.~172--174).

\paragraph{Temperature schedule and logging.} Each goal gets \lean{--rounds} rounds over
every listed model; round 0 samples at temperature $0.2$ and later rounds at $0.8$
(l.~142), so a deterministic-leaning first attempt is followed by more diverse ones. Every
attempt, accepted or not, is appended as one JSON line to
\lean{.leancache/prover\_attempts.jsonl} with the model, round, temperature, latency,
outcome and the first compiler error (ll.~145--170), ``so tier/model effectiveness is
measured, not asserted'' (ll.~11--12).

\paragraph{The compile flags.} The single-file compile passes the two heartbeat and
recursion-depth flags of \cref{tab:pipeline:failures} (ll.~54--57). This line did not
exist at the start of the run; \cref{sec:pipeline:heartbeat} explains what its absence
cost.

\begin{physicsbox}[Why a small local prover is worth running at all]
A 7B-parameter model quantized to 8 bits fits on a 16\,GB GPU and answers in seconds; it
cannot reason about a general $d$-dimensional block matrix, but it has seen a great deal of
Mathlib and reliably produces \lean{decide}, \lean{norm\_num}, \lean{fin\_cases \ldots\ <;>
simp}, and \lean{rfl} proofs for closed computations. The measured median latency was
$4.4$\,s per attempt after a $\sim$170\,s cold load \source{STREAM2\allowbreak\_WORKFLOW.md §6,
l.~111}. Because the kernel checks every candidate, a wrong guess costs a few seconds and
nothing else. This is the sense in which the local tier is ``free'': not that it is
smart, but that being wrong is cheap and being right is verified.
\end{physicsbox}

%=====================================================================================
\section{Routing by the shape of the goal}\label{sec:pipeline:routing}

The cascade of \cref{fig:pipeline:flow} was the plan. What the run measured is that the
split between tiers is ``sharp and predictable'' \source{LL.md §S2.1, l.~24}, and that it
follows the \emph{shape} of the goal rather than its difficulty as a human would judge it.
Three shapes recur.

\begin{description}[leftmargin=*,itemsep=2pt]
\item[Concrete goals.] Fixed matrices, numerals, finite case tables: ``the 64-entry E8
  inverse check'', the signatures of the K3 and Mukai lattices, $\chi(K3\times K3)/24=24$,
  the first five EOT coefficients being $M_{24}$ dimensions. The kernel closes these by
  computation once told which computation to run, and the local prover reliably tells it
  \source{LL.md §S2.1, ll.~25--27}.
\item[Short symbolic algebra.] Statements for general $d$ whose proof is a few rewrites:
  $\eta^2=1$, $\eta\in\Odd{d}$, products of $\Odd{d}$ elements are in $\Odd{d}$, a
  projector is symmetric. Haiku ``closed most'' of these \source{LL.md §S2.1, l.~27}.
\item[Goals needing an idea.] ``Block-matrix identities with inverses,
  positive-definiteness, the diagonal + strict-upper decomposition''---goals where the
  proof requires choosing a decomposition or an auxiliary object that is not in the
  statement. These went to Sonnet \source{LL.md §S2.1, ll.~28--29}.
\end{description}

The improvement list turns the observation into an automatic rule for the next run:
``Concrete goals (numerals, fixed matrices, \lean{fin\_cases}) go to T3
\texttt{prover\_loop.py}. Short symbolic algebra goes to Haiku. Anything with inverses,
positivity, or general $n$ goes straight to Sonnet'' \source{LL.md §S2-I.2, ll.~122--125}.
Note the change from the plan: the plan escalated every goal through every tier; the
measured rule skips tiers for goals whose shape predicts failure, because ``Haiku's success
rate there was low and its reports unreliable'' (l.~125).

\begin{remark}[What the prompts contained]
The comparison between tiers is not a controlled experiment, and the source says so. In
P2.1 Haiku received ``no math hints'', while Sonnet's prompt for
\lean{even\_quadratic\_form\_of\_even\_diag} ``included a strategy hint \ldots\ that the
Haiku run did not get---not a like-for-like comparison'' \source{STREAM2\allowbreak\_WORKFLOW.md §6,
ll.~113--114}. From P2.2 on, ``T2/T1 prompts carried orchestrator-written math sketches''
so the later closure rates ``are therefore \emph{with} sketches''
\source{STREAM2\allowbreak\_WORKFLOW.md §6, ll.~151--152}. A student repeating the experiment should
decide the prompt policy in advance and hold it fixed.
\end{remark}

%=====================================================================================
\section{Measured results}\label{sec:pipeline:results}

This section quotes the pilot numbers exactly as recorded in the workflow document
\source{STREAM2\allowbreak\_WORKFLOW.md §6, ll.~105--196}. They are counts of goals closed at the
kernel gate, not of proofs claimed: ``every agent report was re-compiled and
statement-integrity-checked by the orchestrator before counting'' (ll.~117--119).

\subsection{Phase P2.1: the lattice layer}

The first phase had 26 goals and was the only one run without hints at T2.
\Cref{tab:pipeline:p21} is the record.

\begin{table}[ht]\centering\small
\caption{P2.1 lattice layer, 26 goals, 26/26 closed \source{STREAM2\allowbreak\_WORKFLOW.md §6,
ll.~107--114}. ``Attempted'' counts goals reaching the tier; a goal closed at T3 was not
attempted at T2.}
\label{tab:pipeline:p21}
\begin{tabular}{@{}llccp{6.1cm}@{}}\toprule
Tier & Model & Attempted & Closed & Notes \\ \midrule
T3 & DeepSeek-Prover-V2-7B & 26 & 18 & every concrete goal; 0 of the 8 general-$n$/general-$d$
goals; median latency 4.4\,s/attempt after $\sim$170\,s cold load \\[2pt]
T3 & Goedel-Prover-V2-8B & 8 & 0 & not a capability verdict (\cref{sec:pipeline:thinking});
dropped from routing \\[2pt]
T2 & Haiku (no hints) & 8 & 5 & \lean{isUnimodular\_of\_mul\_eq\_one}, \lean{e8Neg\_unimodular},
\lean{eta\_mul\_self}, \lean{eta\_isODD}, \lean{isODD\_mul} \\[2pt]
T1 & Sonnet & 3 & 3 & \lean{even\_quadratic\_form\_of\_even\_diag},
\lean{thetaShift\_isODD}, \lean{basisChange\_isODD} \\
\bottomrule\end{tabular}
\end{table}

The conclusion drawn at the time: ``the local prover is a very cheap, very reliable closer
for concrete/computational goals and useless for symbolic ones; Haiku closes most short
symbolic algebra; Sonnet was needed only for goals requiring a genuine proof idea. Routing
T3 $\to$ T2 $\to$ T1 put Sonnet on 3 of 26 goals (11.5\%)'' (ll.~121--124). The reader can
check $3/26=0.1154$. Gate G3 for the phase: five modules green, 0 \lean{sorry} warnings,
28 theorems audited, 0 failing (ll.~116--117); the count exceeds 26 because helper lemmas
proved along the way are audited too.

\subsection{Phases P2.3--P2.6, P2.2 and the trace bound}

The second block of work, 23 goals across five modules, is recorded per module
(\cref{tab:pipeline:p23}). Here the split by shape is visible module by module: the
Mukai pairing, the tadpole counts and the EOT decompositions are concrete and went mostly
to T3; the $E_8$ positive-definiteness certificate and the generalized-metric identities
are symbolic with inverses and went almost entirely to T1.

\begin{table}[ht]\centering\small
\caption{P2.3/P2.4-core/P2.5/P2.6, 23 goals, 23/23 closed, by module and closing tier
\source{STREAM2\allowbreak\_WORKFLOW.md §6, ll.~134--143}.}
\label{tab:pipeline:p23}
\begin{tabular}{@{}p{5.4cm}p{1.1cm}p{1.6cm}p{1.6cm}p{1.6cm}@{}}\toprule
Module & Goals & T3 DeepSeek & T2 Haiku & T1 Sonnet \\ \midrule
\texttt{Lattice.Mukai} & 5 & 3 & 2 & --- \\
\texttt{Lattice.E8PosDef} (LDL$^{\mathsf T}$, $\det=1$, positive definite) & 3 & 0 & 0 & 3 \\
\texttt{DFT.GeneralizedMetric} ($(\eta\HH)^2=1$, symmetry, $G\mapsto G^{-1}$, mass covariance) & 6 & 0 & 1 & 5 \\
\texttt{Flux.Tadpole} ($\chi(K3\times K3)/24=24$, $\chi(\KT)=0$, D3 budget) & 5 & 4 & 1 & --- \\
\texttt{Moonshine.EOT} (EOT (1.12), (1.14), (1.15) vs.\ the $M_{24}$ table) & 4 & 2 & 2 & --- \\ \midrule
Total & 23 & 9 & 6 & 8 \\
\bottomrule\end{tabular}
\end{table}

Two smaller packages complete the tally. P2.2, the factorized dualities, had 10 goals
including the helper \lean{proj\_comm}: ``T3 DeepSeek 1 (\lean{factorized\_two}), T2 Haiku 2
kept (\lean{proj\_transpose}, \lean{chargeNorm\_even}; its other claimed solutions did not
compile and were reset), T1 Sonnet 7 (6 goals + helper \lean{proj\_comm}), $\sim$6 compile
runs in under 7 minutes once given the real heartbeat flags and a wall-clock limit''
\source{STREAM2\allowbreak\_WORKFLOW.md §6, ll.~172--177}. The dual-scale trace bound had 6 goals:
``T3 DeepSeek 1 (\lean{dualScale\_eq}), T2 Haiku 3 (\lean{dualScale\_one},
\lean{dualScale\_inv}, \lean{circle\_effective\_scale\_ge\_two}; the agent ran $\sim$6.6\,h
wall-clock and needed a deadline), T1 Sonnet 2 (\lean{dualScale\_ge} via
$\Tr((G-1)G^{-1}(G-1))\ge0$, \lean{dualScale\_circle}) in $\sim$6\,min'' (ll.~179--183).

\subsection{The final gate and the whole-run tally}

The final gate at 2026-09-17 05:07 UTC reports ``roadmap 18/18, designed goals 99/99'':
\texttt{lake build DualScaleStream2} with 3670 jobs and 0 \lean{sorry};
\texttt{tools/axiom\_audit.py} on \lean{DualScaleStream2} with 99 theorems and 0
failing; Stream 1 re-gated in the same session at 3353 jobs, 326 theorems, 0 failing;
statements locked \source{STREAM2\allowbreak\_WORKFLOW.md §6, ll.~185--189}. The last round
of proving is itemized in prose: DeepSeek closed the Mirror identity (3 goals),
\lean{genMetric\_bshift\_cancel}, \lean{mul\_jMat\_mul\_transpose}
and \lean{basisChange\_mul}; Haiku closed one Spectrum lemma, the $\SLZ$ product (2) and 5
of 6 section-condition goals; Sonnet closed reflections (4 + helper), $B$-shift covariance
(2), spectrum equivalence (3 + helper) and flux integrality (4), ``each in under 7
minutes'' (ll.~191--196).

\begin{example}[Adding up the tiers]\label{ex:pipeline:tally}
Summing the four itemized packages gives, per closing tier,
\[
\text{T3: } 18+9+1+1=29,\qquad \text{T2: } 5+6+2+3=16,\qquad \text{T1: } 3+8+7+2=20,
\qquad \text{total } 65 .
\]
The last round adds $6$ (T3), $8$ (T2) and, counting helpers, $4+1+2+1+3+1+4=16$ (T1),
for $95$ attributed closures against $99$ designed goals. The source does not itemize the
difference; the prose records at least one goal ``repaired by the orchestrator'' after a
Haiku proof ``failed to parse (operator precedence)'' (ll.~193--194), and helper lemmas are
counted inconsistently between packages. The student should take from this what the
document itself takes: the per-package tables are exact, the whole-run split by tier is
approximate to within a few goals, and the only number with a gate behind it is 99/99.
\end{example}

\begin{example}[Reading the attempt log]\label{ex:pipeline:log}
The log \texttt{.leancache/prover\_attempts.jsonl} is the primary record of the
T3 tier. On the checkout used for this chapter it holds 155 lines. The author tabulated
them with a ten-line Python script: 151 attempts by DeepSeek-Prover and 4 by
Goedel-Prover; of the DeepSeek attempts 35 were accepted, 59 produced no extractable
proof, and 57 compiled with errors (29 with a single error, 15 with two, and a tail up to
101 errors when a candidate broke the file's structure). The 35 acceptances are 35
distinct theorems, among them \lean{cartanE8Inv\_mul}, \lean{first\_five\_are\_irreps},
\lean{tadpole\_budget}, \lean{factorized\_two} and \lean{dualScale\_eq}.
The median logged latency over all DeepSeek attempts is
$6.9$\,s (maximum $171.5$\,s, the cold load), higher than the $4.4$\,s quoted for P2.1
alone, because later files were longer and the prompt grows with the file. The acceptance
rate per attempt, $35/151\approx23\%$, understates the per-goal rate: the loop keeps trying
a goal it cannot close for every model and round, so failures cluster on the symbolic goals.
The log does not record the T2 and T1 tiers at all, which is a gap the next run should
close (Exercise~\ref{ex:pipeline:logging}).
\end{example}

%=====================================================================================
\section{Failures and the safeguards they produced}\label{sec:pipeline:failures}

Every safeguard in the current pipeline is the residue of a failure that was observed,
written down, and traced to a cause. \Cref{tab:pipeline:failures} is the index; the
subsections give the mechanism of each. The reader who has run any long computation will
recognize the pattern: most of the failures are not failures of mathematics but of
\emph{measurement}---a report that did not match the file, a budget that did not match
the build, a timebox that did not bound time.

\begin{table}[ht]\centering\small
\caption{Recorded failures of the run and the safeguard each produced. W =
\texttt{STREAM2\_WORKFLOW.md}, L = \lean{LL.md}.}
\label{tab:pipeline:failures}
\begin{tabular}{@{}p{4.1cm}p{6.3cm}p{3.6cm}@{}}\toprule
Failure & Safeguard & Source \\ \midrule
Agent reports misstate compile status (``5 clean \lean{sorry}s'' = 8 errors; ``no errors''
= 3; ``0 errors'' with a parse failure) & recompile every report; restore unsolved goals to
exactly \lean{sorry}; prompt line ``final compile 0 errors, paste the literal compiler
output'' & W ll.~153--157, 192--194; L §S2.3 \\[2pt]
Thinking prover spends its whole budget without answering & inspect \texttt{done\_reason} and
the \texttt{thinking} field before judging capability; drop the model at this GPU speed & W
l.~112; L §S2.2 \\[2pt]
\texttt{lake env lean} ignores \texttt{lakefile} options: 200000 heartbeats instead of 1000000
& pass \lean{-DmaxHeartbeats=1000000 -DmaxRecDepth=8000} to every single-file compile & W
ll.~158--163; L §S2.4 \\[2pt]
Hand-entered $M_{24}$ table wrong for months & guard every table with a global consistency
theorem (Burnside) & W ll.~129--132; L §S2.6 \\[2pt]
\lean{native\_decide} in a ``clean'' library & axiom audit at every gate; forbidden-token
filter in the prover loop & W ll.~127--128; L §S2.5 \\[2pt]
Agents run for hours on 1--3 goals & wall-clock limits; at most 2--3 compile-heavy agents
at once & W ll.~164--167; L §S2.8 \\[2pt]
Goal times out on a literal $64$-entry product & re-derive the route (sum of squares) before
escalating the tier & W ll.~168--170; L §S2.7 \\[2pt]
Audit reports ``89 failing (MISSING)'' on a stale build & audit exits 2 with \texttt{AUDIT
ERROR} when the probe does not elaborate & L §S2.11 \\
\bottomrule\end{tabular}
\end{table}

\subsection{Misreported compile status}\label{sec:pipeline:misreport}

The hosted models were asked to prove goals and to report what they had done. Twice in the
second block of work ``Haiku reports misstated compile status: one claimed 5 clean
\lean{sorry}s but left 8 compile errors (\texttt{GeneralizedMetric}); one claimed `no
compilation errors' with 3 errors (\texttt{Factorized}). Both caught only because the
orchestrator recompiles every report (producer $\ne$ verifier)'' \source{STREAM2\allowbreak\_WORKFLOW.md
§6, ll.~153--157}. In the final round a report ``again claimed `0 errors' while one proof
failed to parse (operator precedence)'' (ll.~193--194), and a third report ``said a
definition `could not be found' when it existed in a file this session had written''
\source{LL.md §S2.3, ll.~41--42}.

The consequence for the pipeline is that a report is an \emph{input}, never a result. The
rule now in force has three parts: ``Recompile every report, re-check that every statement
and definition is byte-identical to what was designed, and restore unsolved goals to exactly
\lean{sorry} before escalating'' \source{LL.md §S2.3, ll.~42--44}. The second part is the
statement lock; the third is what makes escalation sound---the next tier must receive the
reviewed statement with an empty proof, not a half-repaired file. The prompt itself now
carries the line ``final compile 0 errors, paste the literal compiler output''
\source{LL.md §S2-I.3, ll.~127--128}, which asks for evidence rather than a verdict.

\subsection{The thinking model that never answered}\label{sec:pipeline:thinking}

The second local prover, Goedel-Prover-V2-8B, closed 0 of 8 goals. The record refuses to
call this a capability verdict: ``with thinking on it spent 8192 tokens / 615\,s without
emitting an answer; with thinking off it drifted off-target and emitted \lean{sorry}
scaffolds'' \source{STREAM2\allowbreak\_WORKFLOW.md §6, l.~112}. The arithmetic explains the first
mode: $8192$ tokens in $615$\,s is about $13$ tokens per second, the speed of an 8B model
on a T4, and the whole generation budget was consumed by the model's private reasoning
before any proof text was produced. The lesson is stated as a diagnostic: ``Check
\texttt{done\_reason} and the separate \texttt{thinking} field before concluding a model can't
prove something. It was a configuration failure at T4 speed, not a capability verdict''
\source{LL.md §S2.2, ll.~36--37}. The four Goedel attempts in the log all carry the reason
\texttt{no\_proof\_extracted}, with a median latency of $97$\,s, consistent with this
reading. The model was ``dropped from routing at T4 speed'', a decision about this hardware,
not about the model.

\subsection{The heartbeat mismatch}\label{sec:pipeline:heartbeat}

Lean bounds the work of elaborating a declaration by a count of ``heartbeats''
(allocation events); a proof that exceeds the bound fails with a timeout. The project's
\texttt{lakefile.lean} raises the bound to $10^6$ from the default $2\times10^5$
(\cref{ch:architecture}). But ``Package \texttt{leanOptions} \ldots\ apply to \lean{lake
build}, not to bare \texttt{lake env lean file.lean}, which uses Lean's default 200000. Agents
and \texttt{tools/prover\_loop.py} were running a stricter gate than the real build, and one
Haiku agent abandoned \lean{e8\_LDL} on a timeout the real budget doesn't hit''
\source{LL.md §S2.4, ll.~48--52}.

The direction of the error is what saved the run. ``A stricter gate can reject valid
proofs; it can never accept invalid ones, so no wrong result entered the repo---it only
cost time'' \source{LL.md §S2.4, ll.~52--54}. The fix is the flag line in
\texttt{prover\_loop.py} quoted in \cref{sec:pipeline:loop} and the same flags in every later
prompt \source{STREAM2\allowbreak\_WORKFLOW.md §6, ll.~161--163}. Notice how the two compiles differ:
the build reads options from the package file, the single-file compile takes them from the
command line, and a harness that uses the second must copy the first by hand. A student
building their own loop should compile one known-good heavy proof through the harness
before trusting its rejections.

\subsection{Wall-clock, not attempts}\label{sec:pipeline:wallclock}

``Several Haiku agents ran for hours on a handful of goals (one: 3 goals, 6.7\,h),
dominated by CPU contention from up to six concurrent \lean{lean} processes on one VM, not
by reasoning. Timeboxes stated as `compile attempts' did not bound wall-clock time''
\source{STREAM2\allowbreak\_WORKFLOW.md §6, ll.~164--167}. The remedy has two parts: a wall-clock
limit per agent ($\approx45$\,min for Haiku, $\approx90$\,min for Sonnet) and a concurrency
budget of ``at most 3 compile-heavy agents'' \source{LL.md §S2-I.3, ll.~126--128}. The
effect was immediate and is recorded in the P2.2 line: about six compile runs ``in under 7
minutes once given the real heartbeat flags and a wall-clock limit---compare hours for the
earlier unbounded Haiku agents'' \source{STREAM2\allowbreak\_WORKFLOW.md §6, ll.~176--177}.

\subsection{The wrong table and the missing axiom check}\label{sec:pipeline:table}

Two defects were found not in Stream 2 but in the Stream 1 library it imports, by gates that
Stream 2 introduced \source{STREAM2\allowbreak\_WORKFLOW.md §6, ll.~126--132}. \Cref{ch:architecture}
tells the story from the gates' side; here is the pipeline's side.

The table \lean{StringDynamics.M24RepDim} of the 26 irreducible dimensions of $\Mtf$
``was wrong: 10395 and 483 duplicated, second 990 and third 1035 missing; squares summed
to 351,061,029 $\ne|M_{24}|$'' (ll.~129--130). The reader can reproduce the wrong sum
without the wrong table: starting from the correct sum $244\,823\,040$ and replacing one
$990$ and one $1035$ by a second $10395$ and a second $483$ changes it by
\[
10395^2+483^2-990^2-1035^2 = 108\,056\,025+233\,289-980\,100-1\,071\,225 = 106\,237\,989,
\]
and $244\,823\,040+106\,237\,989=351\,061\,029$, which is the recorded figure. The table
had passed every build for months because its only theorem ``checked a single entry''
\source{LL.md §S2.6, ll.~64--66}; the guard now in place is the Burnside sum, and the
general rule is ``for every hand-entered table, add the cheapest global consistency theorem
available (sum of squares, total count, known product, symmetry)'' (ll.~66--67). The
second defect, a \lean{native\_decide} in \lean{bps\_ratio\_reduced}, was caught by the
axiom audit and replaced by a kernel \lean{norm\_num} proof (ll.~127--128); the prover loop's
forbidden-token list is the same check applied before compilation rather than after.

\subsection{Route choice beats tier choice}\label{sec:pipeline:route}

The most instructive failure is one where the mathematics, not the tooling, was at fault.
The goal \lean{cartanE8\_posDef}---the $E_8$ Cartan matrix is positive definite---``failed
at T2 through the literal 64-entry $L\cdot D\cdot L^{\mathsf T}$ matrix product. It fell to
T1 once the orchestrator derived (and sympy-verified) the equivalent sum-of-squares identity
$x^{\mathsf T}E_8x=\sum D_ky_k^2$, which \lean{ring} checks directly''
\source{LL.md §S2.7, ll.~70--72}. The rule drawn: ``When a goal times out, look for a
formulation that avoids large definitional unfolding \emph{before} escalating to a costlier
model'' (ll.~72--73).

\begin{example}[The $E_8$ positive-definiteness route]\label{ex:pipeline:e8}
Take the Cartan matrix in the basis used by \lean{cartanE8} \source{E8.lean, ll.~95--103}:
$2$ on the diagonal, $-1$ on the seven bonds of the $E_8$ Dynkin diagram (a chain
$0$--$1$--$2$--$3$--$4$--$5$--$6$ with the eighth node attached to node $4$). Gaussian
elimination without pivoting produces a unit lower-triangular $L$ and a diagonal $D$ with
$LDL^{\mathsf T}=E_8$; the pivots are
\[
D=\diag\Bigl(2,\tfrac32,\tfrac43,\tfrac54,\tfrac65,\tfrac76,\tfrac87,\tfrac18\Bigr),
\]
whose product telescopes to $8\cdot\tfrac18=1=\det E_8$ \source{E8PosDef.lean,
ll.~95--99}. Writing $y=L^{\mathsf T}x$, the quadratic form becomes
\[
x^{\mathsf T}E_8x=(L^{\mathsf T}x)^{\mathsf T}D(L^{\mathsf T}x)=\sum_{k=0}^{7}D_k\,y_k^2,
\qquad y_0=x_0-\tfrac12x_1,\ y_1=x_1-\tfrac23x_2,\ y_2=x_2-\tfrac34x_3,\ \ldots
\]
The author re-derived $L$, $D$ and the identity with sympy: $LDL^{\mathsf T}=E_8$ exactly,
the expanded difference $x^{\mathsf T}E_8x-\sum_kD_ky_k^2$ is the zero polynomial in eight
variables, and the first four $y_k$ agree with the \texttt{set} lines of the Lean proof
\source{E8PosDef.lean, ll.~172--175}. Every $D_k>0$, so the sum is $\ge0$ and vanishes only
if all $y_k=0$; back-substitution through the unit-triangular $L^{\mathsf T}$ then forces
$x=0$. Why did the route matter? The first formulation asked the kernel to multiply three
$8\times8$ matrices symbolically and compare 64 entries \emph{inside} a positivity argument;
the second isolates the 64-entry check as its own concrete lemma \lean{e8\_LDL}
(\lean{fin\_cases} on both indices, then \lean{simp} and \lean{ring}
\source{E8PosDef.lean, ll.~108--112}) and leaves the positivity argument a polynomial
identity that \lean{ring} verifies without unfolding any matrix. That is the concrete/symbolic
split of \cref{sec:pipeline:split} applied inside a single proof.
\end{example}

%=====================================================================================
\section{What the machine has checked}\label{sec:pipeline:machine}

The pipeline itself is not a theorem. The tier assignments, the wall-clock numbers, the
count "18/18 roadmap items, 99/99 goals" are an experimental log, reproducible from
\texttt{docs/STREAM2\allowbreak\_WORKFLOW.md} and \texttt{LL.md} but not kernel-checked; nothing in
this chapter is \TA\ \emph{as a claim about the pipeline}. What is \TA\ are the two
theorems this chapter has followed most closely, one from each end of the tier ladder: the
dual-scale bound, escalated all the way to T1, and the \(E_8\) positivity route of
\cref{ex:pipeline:e8}, whose second formulation is what let T1 close in minutes rather than
time out.

\begin{leanbox}[The dual-scale bound, produced by the process of \cref{sec:pipeline:results}]
\begin{leancode}
noncomputable def dualScale (G : Matrix (Fin d) (Fin d) ℝ) : ℝ :=
  (genMetric G 0).trace
theorem dualScale_ge (G : Matrix (Fin d) (Fin d) ℝ) (hG : G.PosDef) :
    (2 * d : ℝ) ≤ dualScale G := by ...
theorem dualScale_one : dualScale (1 : Matrix (Fin d) (Fin d) ℝ) = 2 * d := by ...
theorem dualScale_circle (R : ℝ) (hR : R ≠ 0) :
    dualScale (!![R ^ 2] : Matrix (Fin 1) (Fin 1) ℝ) = (R + R⁻¹) ^ 2 - 2 := by ...
\end{leancode}
\textbf{Reading.} For every positive-definite $d\times d$ metric $G$, $\Tr G+\Tr G^{-1}\ge
2d$, with equality at the self-dual point $G=1$; on one circle this is $(R+1/R)^2-2$,
recovering \cref{ch:circle}'s $R+1/R\ge2$ after the shift by $-2$. \textbf{Proof idea.}
$\Tr\bigl((G-1)G^{-1}(G-1)\bigr)\ge0$ because it is a trace of the form $M^{\mathsf T}PM$
with $P=G^{-1}$ positive-semidefinite; expanding that same trace with $G^{-1}G=GG^{-1}=1$
gives exactly $\Tr G+\Tr G^{-1}-2d$ \source{TraceBound.lean, ll.~105--143}. No eigenvalue
of $G$ is ever computed, which is why the statement is genuinely $d$-independent and why
T1 closed it in about six minutes once given this route \source{STREAM2\allowbreak\_WORKFLOW.md §6,
final gate}. \textbf{Not proved.} That $G=1$ is the \emph{only} minimizer up to
$O(d)$-conjugation (\cref{ex:open:dualscale-unique} of \cref{ch:open}); the physical
reading of $\mathcal D(G)$ as a minimal length is \TC. \TA
\end{leanbox}

\begin{leanbox}[Route choice, not tier, closed $E_8$ positivity]
\begin{leancode}
theorem e8_LDL : e8L * e8D * e8Lᵀ = cartanE8R := by ...
theorem cartanE8_posDef : cartanE8R.PosDef := by ...
\end{leancode}
\textbf{Reading.} The $E_8$ Cartan matrix is positive definite. \textbf{Proof idea.} An
explicit $LDL^{\mathsf T}$ decomposition with rational pivots turns the quadratic form into
$\sum_k D_k y_k^2$ with every $D_k>0$; \lean{ring} then verifies the polynomial identity
without ever unfolding an $8\times8$ product \source{LL.md §S2.7, ll.~70--72}.
\textbf{Not proved.} That every even unimodular positive- or negative-definite lattice of
rank 8 is isometric to $E_8$ (the classification fact motivating why this matrix was worth
proving positive-definite in the first place) — a literature fact (\TL, Serre, quoted at
\cref{ch:lattices}), not re-derived here. \TA
\end{leanbox}

\begin{center}\small
\begin{tabular}{@{}p{6.3cm}cp{6.5cm}@{}}\toprule
Claim & Tier & Where \\ \midrule
The 99 theorems of \lean{DualScaleStream2}, as stated & \TA & \texttt{lake build}
  3670 jobs, 0 \lean{sorry}; \texttt{axiom\_audit.py} 99/99, axioms \lean{propext},
  \lean{Classical.choice}, \lean{Quot.sound} only \\
"18/18 roadmap items, 99/99 designed goals" & --- & an experimental tally, reproducible
  from \texttt{docs/STREAM2\allowbreak\_WORKFLOW.md}, not a theorem \\
Per-tier closure rates and wall-clock hours & --- & measured this run only; a different
  goal mix or model version would change them \\
That route choice matters more than tier choice & \TC & an engineering judgement of this
  programme, argued from the single case of \cref{ex:pipeline:e8} \\
\bottomrule\end{tabular}
\end{center}

\begin{summarybox}
\begin{itemize}[leftmargin=*,itemsep=1pt]
\item Four tiers, escalated only on failure: T3 a local 7B prover (concrete, closed goals),
  T2 a cheap hosted model (short symbolic algebra), T1 a capable hosted model (anything
  with inverses, positivity, or a free variable $n$), T0 a human or orchestrator review.
\item Design the statement, pre-verify its convention in sympy, and lock it
  (\cref{sec:pipeline:sympy}) \emph{before} any proof search begins; this run interleaved
  the two and paid for it in re-locks (LL.md §S2-I item 1).
\item Never trust a report of "0 errors" or "solved": two Haiku reports this run misstated
  compile status, caught only because the orchestrator recompiled every claim
  (\cref{sec:pipeline:misreport}) --- producer $\ne$ verifier.
\item \texttt{lake env lean} silently uses a stricter heartbeat budget than \texttt{lake build};
  an agent can correctly abandon a goal that the real gate would have accepted
  (\cref{sec:pipeline:heartbeat}).
\item Bound agents by wall-clock, not attempt count: unbounded Haiku agents ran for hours
  on goals Sonnet closed in minutes once the route was right
  (\cref{sec:pipeline:wallclock,sec:pipeline:route}).
\item A green build and zero \lean{sorry} are not the whole gate: axiom audit and statement
  lock catch what a naive grep misses (\cref{sec:pipeline:table}).
\end{itemize}
\end{summarybox}

\section*{Exercises}

\begin{exercise}[$\star$]
State, in one sentence each, the four tiers of \cref{sec:pipeline:tiers} and the one
condition that sends a goal to the next tier. Which tier would you assign
\lean{hyperbolicU\_unimodular} (\cref{ch:lattices}), and which would you assign a general
proof that $\Odd{d}$ is generated by reflections and $\SLZ$ for every $d$?
\end{exercise}

\begin{exercise}[$\star$]
Re-derive by hand, or with sympy, the eight $LDL^{\mathsf T}$ pivots of
\cref{ex:pipeline:e8} from the $E_8$ Cartan matrix of \cref{ch:lattices}, and check that
their product equals $\det E_8=1$.
\end{exercise}

\begin{exercise}[$\star$]
Read \texttt{docs/STREAM2\allowbreak\_WORKFLOW.md §6} and reproduce the final tally --- 99 theorems,
0 failing axiom audit --- against \texttt{papers/book/generated/atlas.md}'s
per-library count for \lean{DualScaleStream2}. Do the two independent counts agree?
\end{exercise}

\begin{exercise}[$\star\star$]
\cref{sec:pipeline:heartbeat} reports that \texttt{lake env lean} ignores
\texttt{maxHeartbeats} set in \texttt{lakefile.lean}. Run a single-file compile of a slow
goal both with and without the flags \lean{-DmaxHeartbeats=\allowbreak 1000000
-DmaxRecDepth=\allowbreak 8000} and record the two heartbeat counts Lean reports on
timeout. Why does this asymmetry make a \emph{false abandon} more likely than a
\emph{false accept}?
\end{exercise}

\begin{exercise}[$\star\star$]
The $E_8$ example replaced a 64-entry matrix product inside a positivity argument by a
polynomial identity in eight free variables. Try the naive route yourself: state
\lean{cartanE8\_posDef} and attempt it with \lean{decide} after unfolding
\lean{Matrix.PosDef} directly (no $LDL^{\mathsf T}$). Where, concretely, does the naive
proof stall?
\end{exercise}

\begin{exercise}[$\star\star$]
\cref{sec:pipeline:misreport} describes a Haiku report that claimed "no compilation
errors" while three errors remained. Propose one line to add to every proving prompt that
would make such a report harder to produce by accident (not by the agent lying, but by an
ambiguous instruction), and explain the failure mode it targets.
\end{exercise}

\begin{exercise}[$\star\star$]\label{ex:pipeline:iff}
The pitfall box of \cref{sec:pipeline:design} records that this run caught a docstring
that had drifted from "if" to "iff" \source{LL.md §S2.13, ll.~109--112}. Pick three
theorems cited in this book whose plain-language "Reading" states an implication, grep
their statements for a genuine \lean{Iff} (or a proof structured as \lean{constructor}
into two directions), and confirm the prose does not silently claim more than the
statement. Report any mismatch you find, however small, as precisely as the pitfall box
does.
\end{exercise}

\begin{exercise}[$\star\star$]\label{ex:pipeline:logging}
\cref{sec:pipeline:results} notes that the local prover's log recorded every T3 attempt
but the run's T2 and T1 attempts were not logged in the same structured way, a gap the
text flags as one the next run should close. Extend \texttt{tools/prover\_loop.py} (or
write a thin wrapper around a hosted-model call) so that every attempt at every tier
writes one line with at least: goal name, tier, model, wall-clock time, and kernel
verdict. Re-derive the P2.1 table's closure counts (\cref{sec:pipeline:results}) from your
log format alone, without consulting the chapter text, and check they agree.
\end{exercise}

\begin{exercise}[$\star\star\star$]
Using \texttt{tools/prover\_loop.py} and a small goal set of your own (five to ten
concrete \lean{decide}-able lemmas), reproduce the T3 tier of this pipeline against a
local prover model. Record closure rate and wall-clock time, and compare them to the
P2.1--P2.6 rows of \cref{sec:pipeline:results}. Does a different goal mix change which
tier dominates?
\end{exercise}

\begin{exercise}[$\star\star\star$]
Design one new statement for a fact not yet in \lean{DualScaleStream2} (see
\cref{ch:open} for candidates), pre-verify its convention symbolically, lock it with
\texttt{tools/statement\_\allowbreak lock.py}, and route it through the tiers of this
chapter by hand. Report which tier closed it and whether the routing rule of
\cref{sec:pipeline:routing} predicted that correctly in advance.
\end{exercise}

\section*{Further reading}

The full log is the workflow document cited throughout this chapter as
\texttt{STREAM2\allowbreak\_WORKFLOW.md}. The lessons distilled from it are
\texttt{LL.md §S2.1--§S2.13} (this run) and
\texttt{§S2-I} (what to change next, read in full at \cref{ch:open}). The tools
referenced in this chapter:
\texttt{tools/prover\_loop.py} (T3), \texttt{tools/axiom\_audit.py} (the axiom gate),
\texttt{tools/statement\_lock.py} (the comment-insensitive lock), and
\texttt{tools/gate.sh} proposed in \texttt{LL.md §S2-I} item 4 to wrap all of them in one
script. For the mathematics produced by this process: \cref{ch:lattices,ch:dualscale} and
the final tally of \cref{sec:pipeline:results}. For what the atlas found once the library
was built --- including the duplication this pipeline's problem-by-problem structure
tends to produce --- \cref{ch:atlas}.
